\documentclass[12pt]{article}
\usepackage[utf8]{inputenc}
\usepackage[T1]{fontenc}
\usepackage{hyperref}
\title{Wintu Language Resources for K-12 Curriculum Development}
\author{Compiled by Claude Sonnet 5.5 (claude-sonnet-5-5) for Andr\'as Kornai}
\date{}
\begin{document}
\maketitle

\section{Introduction}
Wintu (ISO 639-3 wnw) is a Wintuan language of northern California with no first-language speakers in the 21st century (Golla 2011). Revitalization efforts by tribal members are under way.

\section{Primary Educational Resources}
The search results do not show any published K-12 curriculum. The Winnemem Wintu Language Project has been working since 2011 with linguist Stefan Liedtke on an alphabet and learning materials, including a talking dictionary. I saw no URL for the project itself, so it is mentioned here only.

\section{Linguistic and Reference Materials}
\begin{itemize}
\item Harvey Pitkin, \emph{Wintu Dictionary} (1985, University of California Press), a Wintu--English dictionary with an English finderlist. A digital version is on the UBC wiki: \url{https://wiki.ubc.ca/Documentation:RelLex/Wintu_Dictionary_(1985)}
\item Wikipedia, Wintu language: \url{https://en.wikipedia.org/wiki/Wintu_language}
\item Omniglot, Wintu: \url{https://omniglot.com/writing/wintu.htm}
\item UC Berkeley Survey of California and Other Indian Languages, Wintu page: \url{https://cla.berkeley.edu/languages/wintu.html}
\end{itemize}
Alice Schlichter's 1981 Wintu dictionary appears in the Berkeley bibliography, but I saw no URL for it.

\section{Religious and Community Resources}
The Winnemem Wintu Language Project and the Wintu Tribe of Northern California are active in revitalization, according to search summaries. I saw no URLs for their own sites.

\section{Digital and Audio Resources}
\begin{itemize}
\item Endangered Languages Project, Wintu resources: \url{https://endangeredlanguages.com/node/37667/resources}
\item UC Berkeley Language Revitalization listing: \url{https://lx.berkeley.edu/person-type/language-revitalization}
\item Anna Bj\"orklund, UC Berkeley, whose research uses archival recordings for phonetic analysis and revitalization tools (a Wintu linguistics consultant is noted for the Wintu Tribe of Northern California): \url{https://lx.berkeley.edu/people/anna-bj%C3%B6rklund}
\end{itemize}
I found no dedicated Wintu audio archive.

\section{Curriculum Development Context}
With no first-language speakers, a curriculum would be a heritage and revival one, built on Pitkin's dictionary and archival materials. It should be developed with the tribal communities and their language programmes.
\end{document}